\documentclass[12pt,a4paper]{article}

\usepackage[utf8]{inputenc}
\usepackage[T1]{fontenc}
\usepackage[margin=2.5cm]{geometry}
\usepackage{hyperref}
\usepackage{microtype}
\usepackage{parskip}
\usepackage{lmodern}

\hypersetup{colorlinks=true, urlcolor=blue!60!black}

\pagestyle{empty}

\begin{document}

\begin{flushright}
Kundai Farai Sachikonye \\
Auf der Lichtung 19, 82178 Puchheim, Germany \\
\href{https://github.com/fullscreen-triangle}{github.com/fullscreen-triangle} \\
\texttt{kundai.sachikonye@bitspark.com} \\
(+49) 1626386344 \\[6pt]
18 May 2026
\end{flushright}

\vspace{1em}

Institute of Neuroscience and Medicine --- Computational Biomedicine (INM-9) \\
Forschungszentrum Jülich GmbH \\
52428 Jülich, Germany \\
Reference: 2025D-126

\vspace{1.5em}

\textbf{Re: PhD Position --- Simulations and Machine Learning Applied to Proteins (2025D-126)}

\vspace{1em}

Dear Hiring Committee,

I am applying for the doctoral position in simulations and machine learning applied
to proteins at INM-9 (Ref.\ 2025D-126). My background sits at the intersection
that makes this position attractive: I have developed formal ML frameworks for
physics-based simulation surrogates and first-principles computational tools
for molecular structure and dynamics, without having run production MD campaigns
in GROMACS or AMBER. I am therefore a candidate whose ML and molecular physics
depth is established, and for whom the classical MD and enhanced sampling methodology
is what the PhD provides --- the same way experimental microscopy skills are what
a super-resolution physics PhD provides.

\textbf{Machine learning for molecular simulation: surrogate models and training theory.}
The most technically demanding part of ML-augmented MD is not training a neural
network on a trajectory --- it is deciding \emph{what} to train, on \emph{which}
data, and guaranteeing that the surrogate is competent outside the training distribution.
My \href{https://github.com/fullscreen-triangle/purpose}{Purpose} framework addresses
this directly. The \textbf{Backward Training Principle} proves that a surrogate trained
exclusively at the \emph{pure-intent regime} --- the fully-resolved, single-objective
instance of the learning problem --- is $\epsilon$-competent on every constrained
sub-regime by feasibility inclusion, not generalisation. A Forward Training Collapse
Theorem establishes a provably positive Vapnik--Chervonenkis failure gap for the
conventional easy-to-hard approach; a Sample Complexity Theorem gives $(N{+}1)\times$
sample savings for a cascade of depth $N$.
For protein simulation, the pure-intent regime is the fully-folded, minimum free-energy
conformation in the absence of solvent and finite-size artefacts: the regime from which
all realistic simulation conditions are feasibility restrictions. Training an ML
surrogate there, and projecting competence downward to partial folding states and
aqueous conditions, eliminates the need to sample the full configurational space.
The \href{https://github.com/fullscreen-triangle/homo-veloce}{homo-veloce} system
demonstrates the pattern in practice: Gaussian process and graph neural network
physics solvers distilled into lightweight 1D CNNs, LSTMs, and Transformers for
real-time inference --- the same pattern the field uses to replace force evaluations
with learned potentials.

\textbf{First-principles molecular physics: zero-parameter spectroscopy and force fields.}
The \href{https://honjo-masamune.vercel.app/tools}{Borgia} framework derives
electronic transitions, absorption, emission, and vibrational modes of molecules
from bounded phase space geometry alone, with zero empirical parameters.
The Triple Equivalence Theorem establishes that oscillatory dynamics, categorical
state enumeration, and partition operations are formally identical representations
of the same molecular object --- meaning a GPU shader computing the partition
function is formally equivalent to a physical spectrometer.
Validated against NIST for nine elements and six molecules with self-consistency
errors below 1.63\,\%.
For protein ML, this matters because it provides a parameter-free baseline for
molecular force contributions: instead of fitting empirical parameters to reproduce
experimental observables, the observable is derived from geometry.
The same framework --- Beer-Lambert absorption as a Kirchhoff circuit law
(Resonant Stack theorem, transmittance $T = 1/e = 0.3679$ in both representations) ---
applies to light-triggered protein dynamics and photoresponsive residue behaviour,
directly relevant to optogenetics applications in the neuroscience context of INM-9.

\textbf{Protein state and trajectory analysis.}
The \href{https://shakespear-nine.vercel.app/instrument}{Shakespear} platform
provides protein folding, trajectory, catalysis, and dynamics analysis via
Kuramoto oscillator synchronization in S-entropy coordinate space.
The S-entropy coordinates $(S_k, S_t, S_e) \in [0,1]^3$ --- knowledge entropy
(spectral power distribution), temporal entropy (frequency span), and evolution
entropy (coherence dimensionality) --- provide a continuous, reparameterisation-invariant
embedding of protein conformational states.
Trajectory completion via the fixed-point operator $T = P \circ B \circ K$
(Kirchhoff propagation, backward Viterbi inference, thermodynamic projection)
reconstructs full conformational trajectories from sparse observations, with Banach
contraction guaranteeing unique convergence.
The Triple Equivalence gives the biological interpretation: every stable protein
conformation is a bounded oscillatory system, and the S-entropy coordinates are
its complete characterisation at the level of classical statistical mechanics.
This is a complementary tool to conventional MD: where MD integrates Newton's
equations forward, the fixed-point operator completes trajectories backward from
known endpoint states --- applicable to folding pathway reconstruction and
transition-state characterisation without forward simulation.

\textbf{HPC computing and large-scale data.}
I have practical experience working with large-scale scientific datasets on
distributed computing infrastructure. The \href{https://github.com/fullscreen-triangle/lavoisier}{Lavoisier}
pipeline processes mass spectrometry datasets exceeding 100\,GB using Ray
distributed computing, memory-mapped I/O, and HPC-scale parallelism; I have
written the associated compute allocation documentation for both local cluster
and federated multi-site configurations. The experience of structuring a
computational proposal --- specifying system sizes, simulation timescales, memory
requirements, and expected output volume --- transfers directly to the GPU/CPU
allocation proposals for JUWELS and the FZ Jülich HPC infrastructure.

\textbf{Background.}
I hold an MSc in Computational Biology (Universität Tübingen) and an MSc in
Pharmaceutical Biotechnology (HAW Hamburg), and I completed doctoral research
in Computational Lipidomics at TUM. I have not run production MD campaigns ---
I am aware the PhD provides that training --- but the molecular physics understanding
and the ML methodology I bring are at the level the position requires.
I am legally resident in Germany with full work authorisation.
English is native; German is conversational.
I am prepared to relocate to Jülich.

\vspace{1.5em}
Yours sincerely,

\vspace{2.5em}
\textbf{Kundai Farai Sachikonye}

\end{document}
